\chapter{Introduction}
\label{chap:Introduction}

%---------------------------- SECTION ----------------------------
\section{A Sister Volume to \textit{Prudentia}}

\paragraph{This book is one of a planned series. The base volume, \textit{Prudentia}, sets out the profession-agnostic method. Each subsequent volume --- of which this is the first --- applies that method to a single profession. Future volumes (\textit{Prudentia Medica}, \textit{Prudentia Iuridica}, and others) will follow the same scaffold.}

\paragraph{\textit{Prudentia Veterinaria} is the veterinary volume in the \textit{Prudentia} series.}

\paragraph{Veterinary practice is unusually exposed to multi-axis risk. A single event can be clinical and regulatory at once; a medicines issue can become both a welfare issue and a criminal law issue; a staffing gap can shift from operational inconvenience to patient harm within a single shift. You are therefore rarely assessing one thing in isolation. You are assessing systems: people, process, environment, equipment, records, and oversight.}

\paragraph{This book assumes that your goal is not merely to pass inspection, but to run a safer and more reliable practice. Compliance remains essential, but compliance is treated as a floor rather than a ceiling. The chapters that follow are designed to support decisions in real settings: consult rooms, theatres, kennels, imaging suites, farm visits, dispensaries, and out-of-hours rotas.}

%---------------------------- SECTION ----------------------------
\section{Who this book is for}

\paragraph{\textit{Prudentia Veterinaria} is written for practice principals, clinical directors, head nurses, practice managers, and RCVS PSS-focused governance leads who carry accountability for risk decisions. It is also designed for locums and temporary leaders who need to understand local risk architecture quickly and safely.}

\paragraph{The intended reader is not only a compliance specialist. It includes clinicians who own real-time decisions, managers who assign resources, and assessors who evaluate whether governance is genuinely operating. The chapters are therefore practical, structured, and written to support implementation and assurance.}

%---------------------------- SECTION ----------------------------
\section{Methodology used in this book}

\paragraph{This volume applies the \textit{Prudentia} five-step method directly to veterinary contexts: identify hazards, determine who/what is affected, evaluate risk, select controls, and record/review. The value lies in repeatability and comparability across domains.}

\paragraph{Every chapter uses the same structure: framing, legal/professional framework, five-step application, hierarchy of controls, cadence and triggers, common failures, and summary tools. This repeating rhythm supports delegation, training, and governance oversight.}

\barquote{A method is only useful when different teams can apply it in different settings and still produce defensible decisions.}

%---------------------------- SECTION ----------------------------
\section{Legal context and devolved variation}

\paragraph{\textit{Prudentia Veterinaria} is UK-centred. It references the Veterinary Surgeons Act 1966, the RCVS Code, RCVS PSS standards, the Veterinary Medicines Regulations, health and safety law, COSHH, radiation controls, fire safety law, and data protection law.}

\paragraph{Material variation exists across the four UK nations in some statutes and operational pathways. Animal welfare legislation and fire safety frameworks differ by jurisdiction. Practices working across borders should maintain local legal mapping and update assessments as guidance evolves.}

\paragraph{Where uncertainty exists on a precise clause number, this volume states obligations in practical terms rather than fabricating citation detail. That approach supports accuracy and professional duty under the RCVS Code to remain current.}

%---------------------------- SECTION ----------------------------
\section{How to use Prudentia Veterinaria in governance}

\paragraph{Use this book as a build tool, an audit benchmark, and a training framework. Build chapter-aligned assessments for depth, then aggregate residual risk and action status into a practice risk register for oversight.}

\paragraph{For inspection readiness, treat documentation as evidence of operational reality. Include ownership, competency records, maintenance records, incident learning, and review triggers. A polished document without evidence of control operation is weak assurance.}

\paragraph{For continuous improvement, apply scheduled review at the chapter minimum and tighten cadence where risk remains high or operating conditions change. Integrate chapter outputs with SEA, clinical audit, and PSS quality improvement cycles.}

\barquote{In Chapter~\ref{chap:What is Risk in a Veterinary Context?}, we begin by defining risk in specifically veterinary terms and setting the language for the whole volume.}

